% GENERATED FILE. Edit canonical lessons, then run npm run generate:books.
% canonical-source-hash: fnv1a64:9a55fe4f3ae536fc
% canonical-lessons: GE-C112-teppich, GE-C112-fernseher, GE-C112-computer, GE-C112-kuehlschrank, GE-C112-herd

\chapter{Rug, Television, Computer, Fridge, Cooker}
\label{ch:ge-rug-television-computer-fridge-cooker}
\hlchaptermodality{112}

\begin{chapteropening}
\textbf{By the end of this chapter:} \emph{I can say a rug, a television, a computer, a fridge and a cooker.}

Complete the last lesson of chapter 112: I can say a rug, a television, a computer, a fridge and a cooker.
\end{chapteropening}

\section[der Teppich]{der Teppich — a rug, a carpet}

\label{lesson:GE-C112-teppich}



\begin{quote}
Before the new one: say the German for a lift, then the German for to fly.

\emph{Read it:} \emph{der Flughafen, das Hotel, das Kino} aloud as one line, without stopping between the words
\end{quote}

\subsection*{You'll want to know: der Teppich}
\textbf{der Teppich} — \textquotedblleft{}a rug, a carpet\textquotedblright{}.

\begin{cousinweb}
\textbf{Teppich} $\leftarrow$ Old High German \emph{teppih}, from Latin \emph{tapētum}, \textquotedblleft{}carpet, tapestry\textquotedblright{} (EtymDB 2.1 has the Latin word), itself from Greek — borrowed. English \emph{tapestry} shares the word.
\end{cousinweb}

\begin{grammarlens}[title={What you've built}]
One more word for this chapter.
\end{grammarlens}

\subsection*{Your turn}
\begin{itemize}
\raggedright
  \item \emph{Say it:} \emph{der Teppich}
  \item \emph{Say it:} \emph{der Teppich}, once more
  \item \emph{Say it:} say \emph{der Teppich}
  \item \emph{Recall:} say the German for a lift, then the German for to fly, then say \emph{der Teppich} again
\end{itemize}

\begin{tcolorbox}[breakable,colback=teal!4,colframe=teal!35!black,title={Before you move on}]
What does der Teppich mean? (\textquotedblleft{}A rug, a carpet\textquotedblright{}.) Say it once more.
\end{tcolorbox}
\section[der Fernseher]{der Fernseher — a television}

\label{lesson:GE-C112-fernseher}



\begin{quote}
Before the new one: say the German for to fly, then the German for a rug.
\end{quote}

\subsection*{You'll want to know: der Fernseher}
\textbf{der Fernseher} — \textquotedblleft{}a television\textquotedblright{}.

\begin{cousinweb}
\textbf{Fernseher} is the word for \textquotedblleft{}far\textquotedblright{} + \emph{sehen}, \textquotedblleft{}to see\textquotedblright{}, + \emph{-er}: \textquotedblleft{}the far-seer\textquotedblright{}. It translates \emph{television} piece by piece, Greek \emph{tēle}, \textquotedblleft{}far\textquotedblright{}, + Latin \emph{vīsiō}, \textquotedblleft{}seeing\textquotedblright{}. \emph{sehen} $\leftarrow$ Proto-Germanic \emph{*sehwaną}. (EtymDB 2.1)
\end{cousinweb}

\begin{grammarlens}[title={What you've built}]
One more word for this chapter.
\end{grammarlens}

\subsection*{Your turn}
\begin{itemize}
\raggedright
  \item \emph{Say it:} \emph{der Fernseher}
  \item \emph{Say it:} \emph{der Fernseher}, once more
  \item \emph{Say it:} say \emph{der Fernseher}
  \item \emph{Recall:} say the German for to fly, then the German for a rug, then say \emph{der Fernseher} again
\end{itemize}

\begin{tcolorbox}[breakable,colback=teal!4,colframe=teal!35!black,title={Before you move on}]
What does der Fernseher mean? (\textquotedblleft{}A television\textquotedblright{}.) Say it once more.
\end{tcolorbox}
\section[der Computer]{der Computer — a computer}

\label{lesson:GE-C112-computer}



\begin{quote}
Before the new one: say the German for a rug, then the German for a television.
\end{quote}

\subsection*{You'll want to know: der Computer}
\textbf{der Computer} — \textquotedblleft{}a computer\textquotedblright{}.

\begin{cousinweb}
\textbf{Computer} $\leftarrow$ English \emph{computer}, from \emph{compute} + \emph{-er} (EtymDB 2.1), and \emph{compute} $\leftarrow$ Latin \emph{computāre}, \textquotedblleft{}to reckon\textquotedblright{} — borrowed.
\end{cousinweb}

\begin{grammarlens}[title={What you've built}]
One more word for this chapter.
\end{grammarlens}

\subsection*{Your turn}
\begin{itemize}
\raggedright
  \item \emph{Say it:} \emph{der Computer}
  \item \emph{Say it:} \emph{der Computer}, once more
  \item \emph{Say it:} say \emph{der Computer}
  \item \emph{Recall:} say the German for a rug, then the German for a television, then say \emph{der Computer} again
\end{itemize}

\begin{tcolorbox}[breakable,colback=teal!4,colframe=teal!35!black,title={Before you move on}]
What does der Computer mean? (\textquotedblleft{}A computer\textquotedblright{}.) Say it once more.
\end{tcolorbox}
\section[der Kühlschrank]{der Kühlschrank — a fridge}

\label{lesson:GE-C112-kuehlschrank}



\begin{quote}
Before the new one: say the German for a television, then the German for a computer.
\end{quote}

\subsection*{You'll want to know: der Kühlschrank}
\textbf{der Kühlschrank} — \textquotedblleft{}a fridge\textquotedblright{}.

\begin{cousinweb}
\textbf{Kühlschrank} is the word for \textquotedblleft{}cool\textquotedblright{} + \emph{Schrank}, \textquotedblleft{}cupboard\textquotedblright{}: a cooling cupboard. The \textquotedblleft{}cool\textquotedblright{} word, Old High German \emph{kuoli}, is listed in EtymDB 2.1 beside English \emph{cool}.
\end{cousinweb}

\begin{grammarlens}[title={What you've built}]
One more word for this chapter.
\end{grammarlens}

\subsection*{Your turn}
\begin{itemize}
\raggedright
  \item \emph{Say it:} \emph{der Kühlschrank}
  \item \emph{Say it:} \emph{der Kühlschrank}, once more
  \item \emph{Say it:} say \emph{der Kühlschrank}
  \item \emph{Recall:} say the German for a television, then the German for a computer, then say \emph{der Kühlschrank} again
\end{itemize}

\begin{tcolorbox}[breakable,colback=teal!4,colframe=teal!35!black,title={Before you move on}]
What does der Kühlschrank mean? (\textquotedblleft{}A fridge\textquotedblright{}.) Say it once more.
\end{tcolorbox}
\section[der Herd]{der Herd — a cooker}

\label{lesson:GE-C112-herd}



\begin{quote}
Before the new one: say the German for a computer, then the German for a fridge.
\end{quote}

\subsection*{You'll want to know: der Herd}
\textbf{der Herd} — \textquotedblleft{}a cooker\textquotedblright{}.

\begin{cousinweb}
\textbf{Herd} $\leftarrow$ Old High German \emph{herd} $\leftarrow$ Proto-Germanic \emph{*herþaz}, \textquotedblleft{}hearth\textquotedblright{} — inherited, the same word as English \emph{hearth}. (EtymDB 2.1) A cooker took over the hearth's job.
\end{cousinweb}

\begin{grammarlens}[title={What you've built}]
One more word for this chapter.
\end{grammarlens}

\subsection*{Your turn}
\begin{itemize}
\raggedright
  \item \emph{Say it:} \emph{der Herd}
  \item \emph{Say it:} \emph{der Herd}, once more
  \item \emph{Say it:} say \emph{der Herd}
  \item \emph{Recall:} say the German for a computer, then the German for a fridge, then say \emph{der Herd} again
\end{itemize}

\begin{tcolorbox}[breakable,colback=teal!4,colframe=teal!35!black,title={Before you move on}]
What does der Herd mean? (\textquotedblleft{}A cooker\textquotedblright{}.) Say it once more.
\end{tcolorbox}
